\section{モデリングの目から見た検定1；二群の平均値の差}
\subsection{授業内容}
\subsubsection{科目の中でのこのコマの位置づけ}
データ生成モデリングの観点を踏まえた上で，検定的アプローチとベイズ的アプローチの違いを学ぶ。

二群の平均値の差の検定，いわゆる対応のないt検定の場合は，同一の正規分布から得られたデータに対して平均値の差があると判断して良いかどうかを判断するという枠組みであった。
これらの前提と判断基準を確認し，それがデータ生成モデルの観点ではどのように表されるかを検証する。
ここで結果が分布として推定されること，差があるかないかというのは二群の推定された平均値の差であることから，生成量を使って平均値の差を出力することを考える。ここで効果量に改めて目を向けるとその理解が進む。

\subsubsection{コマ主題細目}
\begin{description}
\item[t検定の仮定] 二群の平均値の差を検定するときはt検定が利用されるが，データが正規分布から得られているという仮定，分散が同じであるという仮定などを踏まえて設計図を書き，これをStanで表現することを考える。あらためてt検定のやり方や結果と比べてみることで，モデリングがデータ生成メカニズムに注目していること，パラメータの推定を行なっていることなどが確認できるだろう。また等分散性の仮定を外す方法についてもすぐに応用ができる。
\begin{flushright}
    帰無仮説検定と二群の差の検定について，$\to$\citet{Yamada2004}や一年次の資料をもとに確認しておく。
\end{flushright}
\item[差の分布] 検定は推測に加えて判断を行なっていた，ということを改めて確認するとともに，帰無仮説検定では母平均の差をターゲットにしていたことを確認する。MCMCは母集団からの代表値であるので，推定された結果を使って差を表現することができる。これはR側で得られたサンプルで行っても良いし，Stanの生成量を使っても良い。ここで\texttt{generated quantities}ブロックの考え方を導入し，平均値の差の分布を確認すること，帰無仮説検定が差の分布の一点についての仮説であったことを確認する。また一方が他方よりも大きくなる確率はどれぐらいかとか，一方と他方が$c$以上に違っている確率はどれぐらいか，と言ったことが生成量を使って計算することができるようともいえる。
\item[帰無仮説検定を省みる] ここまでくると，帰無仮説検定のロジックや考え方について別の視点から見ることができるようになるであろう。まずは帰無仮説と対立仮説という対立のさせ方の不平等さである。帰無仮説は一点についての仮説であり，対立仮説はそれ以外であればなんでも良い，という非対称な関係になっていた。それを省みると，差があるかないかといった二値判断に陥ることがいかに危険であるかがわかるだろう。また量的な判断ができないことから，効果量を合わせて報告することが望ましいとされている。効果量とは，標準化された差の大きさのことであり，生成量を使って簡単に算出することができる。また方向性を持った検定について，片側・両側検定などで考えられてきたが，生成料を使えば自然にそれが検証できることがわかる。ただしこれらの検証の仕方は，今回のデータと仮定されたモデルという前提の上で成立する程度であって，過度な一般化にはならないように注意する必要がある。
\end{description}
\subsubsection{キーワード}
\begin{itemize}
    \item t検定
    \item 生成量
    % \item パラメータ・リカバリ
    % \item 事後予測分布
    \item 効果量
    \item 片側検定，両側検定
\end{itemize}
\subsection{授業情報}
\paragraph{コマの展開方法}
講義/遠隔/演習
\subsubsection{予習・復習課題}
\paragraph{予習}Stanの基本的なブロック構成，設計図からコードに落とすやり方を確認しておこう。
\paragraph{復習}データのサイズが変わるだけでなく，平均値の差，効果量が変化したときのt検定の結果とベイズ推定の結果がどのように変わるのか，様々なケースを想定して「遊んで」みるとよい。加えて，そのほかの仮設検定がどのようなデータ生成メカニズムで表現できるかを考えることは，次回以降の準備にもつながる。